% Contents from body/page-473.tex.
% Original printed page 459.
\BookMark{0}{Bibliography}{p-473}
\ChapterSheet[10.35]{473}{459}{}{%
\vspace*{24pt}%
\noindent{\Large\bfseries Bibliography}\par\vspace{4pt}\hrule
\vspace{169pt}%
\fontsize{9.25}{9.8}\selectfont%
\setlength{\parskip}{0pt}\setlength{\parindent}{0pt}%
\setlength{\emergencystretch}{1.7em}%
\par\vspace{2.15pt}\noindent \hypertarget{bib-473-0}{}Altman, A. and Kleiman, S.\par
\noindent\hangindent=25bp\hangafter=1\hspace*{6bp}%
\makebox[19bp][r]{\hypertarget{bib-473-0-item-1}{}1.\enspace}\emph{Introduction to Grothendieck Duality Theory}, Lecture Notes in Math. \emph{146} Springer-Verlag, Heidelberg (1970), 185 pp.\par
\par\vspace{2.15pt}\noindent \hypertarget{bib-473-1}{}Artin, M.\par
\noindent\hangindent=25bp\hangafter=1\hspace*{6bp}%
\makebox[19bp][r]{\hypertarget{bib-473-2}{}1.\enspace}\emph{Grothendieck topologies}, Harvard Math. Dept. Lecture Notes (1962).\par
\noindent\hangindent=25bp\hangafter=1\hspace*{6bp}%
\makebox[19bp][r]{\hypertarget{bib-artin-2}{}2.\enspace}The implicit function theorem in algebraic geometry, in \emph{Algebraic Geometry, Bombay 1968}, Oxford Univ. Press, Oxford (1969), 13--34.\par
\noindent\hangindent=25bp\hangafter=1\hspace*{6bp}%
\makebox[19bp][r]{\hypertarget{bib-473-3}{}3.\enspace}Some numerical criteria for contractibility of curves on algebraic surfaces, \emph{Amer. J. Math. 84} (1962) 485--496.\par
\noindent\hangindent=25bp\hangafter=1\hspace*{6bp}%
\makebox[19bp][r]{\hypertarget{bib-473-4}{}4.\enspace}Algebraization of formal moduli II: Existence of modifications, \emph{Annals of Math. 91} (1970) 88--135.\par
\par\vspace{2.15pt}\noindent \hypertarget{bib-473-5}{}Atiyah, M. F.\par
\noindent\hangindent=25bp\hangafter=1\hspace*{6bp}%
\makebox[19bp][r]{\hypertarget{bib-atiyah-1}{}1.\enspace}Vector bundles over an elliptic curve, \emph{Proc. Lond. Math. Soc.} (3) VII \emph{27} (1957), 414--452.\par
\par\vspace{2.15pt}\noindent \hypertarget{bib-473-6}{}Atiyah, M. F. and Macdonald, I. G.\par
\noindent\hangindent=25bp\hangafter=1\hspace*{6bp}%
\makebox[19bp][r]{\hypertarget{bib-atiyah-macdonald-1}{}\hypertarget{legacy-b-am}{}1.\enspace}\emph{Introduction to Commutative Algebra}. Addison-Wesley, Reading, Mass. (1969), ix + 128 pp.\par
\par\vspace{2.15pt}\noindent \hypertarget{bib-473-7}{}Baily, W. L., Jr.\par
\noindent\hangindent=25bp\hangafter=1\hspace*{6bp}%
\makebox[19bp][r]{\hypertarget{bib-473-7-item-1}{}1.\enspace}On the automorphism group of a generic curve of genus $>2$, \emph{J. Math. Kyoto Univ. 1} (1961/2) 101--108; correction p.~325.\par
\par\vspace{2.15pt}\noindent \hypertarget{bib-473-8}{}B\u{a}nic\u{a}, C. and St\u{a}n\u{a}\c{s}il\u{a}, O.\par
\noindent\hangindent=25bp\hangafter=1\hspace*{6bp}%
\makebox[19bp][r]{\hypertarget{bib-banica-stanasila-1}{}1.\enspace}\emph{Algebraic Methods in the Global Theory of Complex Spaces}, John Wiley, New York (1976) 296 pp.\par
\par\vspace{2.15pt}\noindent \hypertarget{bib-473-9}{}Berthelot, P.\par
\noindent\hangindent=25bp\hangafter=1\hspace*{6bp}%
\makebox[19bp][r]{\hypertarget{bib-berthelot-1}{}1.\enspace}\emph{Cohomologie Cristalline des Sch\'emas de Caract\'eristique $p>0$}, Lecture Notes in Math. \emph{407}, Springer-Verlag, Heidelberg (1974), 604 pp.\par
\par\vspace{2.15pt}\noindent \hypertarget{bib-473-10}{}Birch, B. J. and Kuyk, W., ed.\par
\noindent\hangindent=25bp\hangafter=1\hspace*{6bp}%
\makebox[19bp][r]{\hypertarget{bib-birch-kuyk-1}{}1.\enspace}\emph{Modular Functions of One Variable, IV (Antwerp)}, Lecture Notes in Math. \emph{476}, Springer-Verlag, Heidelberg (1975).\par
}
% Contents from body/page-474.tex.
% Original printed page 460.
\ChapterSheet[10.25]{474}{460}{\Goto{p-473}{Bibliography}}{%
\fontsize{9.25}{9.8}\selectfont%
\setlength{\parskip}{0pt}\setlength{\parindent}{0pt}%
\setlength{\emergencystretch}{1.7em}%
\par\vspace{2.15pt}\noindent \hypertarget{bib-474-0}{}Bombieri, E.\par
\noindent\hangindent=25bp\hangafter=1\hspace*{6bp}%
\makebox[19bp][r]{\hypertarget{bib-bombieri-1}{}1.\enspace}Counting points on curves over finite fields (d'apr\`es S.~A. Stepanov), \emph{S\'eminaire Bourbaki 430} (1972/73).\par
\par\vspace{2.15pt}\noindent \hypertarget{bib-474-1}{}Bombieri, E. and Husemoller, D.\par
\noindent\hangindent=25bp\hangafter=1\hspace*{6bp}%
\makebox[19bp][r]{\hypertarget{bib-bombieri-husemoller-1}{}1.\enspace}Classification and embeddings of surfaces, in \emph{Algebraic Geometry, Arcata 1974}, Amer. Math. Soc. Proc. Symp. Pure Math. \emph{29} (1975), 329--420.\par
\par\vspace{2.15pt}\noindent \hypertarget{bib-474-2}{}Borel, A. and Serre, J.-P.\par
\noindent\hangindent=25bp\hangafter=1\hspace*{6bp}%
\makebox[19bp][r]{\hypertarget{bib-borel-serre-1}{}1.\enspace}Le th\'eor\`eme de Riemann-Roch, \emph{Bull. Soc. Math. de France 86} (1958), 97--136.\par
\par\vspace{2.15pt}\noindent \hypertarget{bib-474-3}{}Borelli, M.\par
\noindent\hangindent=25bp\hangafter=1\hspace*{6bp}%
\makebox[19bp][r]{\hypertarget{bib-474-3-item-1}{}1.\enspace}Divisorial varieties, \emph{Pacific J. Math. 13} (1963), 375--388.\par
\par\vspace{2.15pt}\noindent \hypertarget{bib-474-4}{}\hypertarget{legacy-b-bour}{}Bourbaki, N.\par
\noindent\hangindent=25bp\hangafter=1\hspace*{6bp}%
\makebox[19bp][r]{\hypertarget{bib-474-4-item-1}{}1.\enspace}\emph{Alg\`ebre Commutative}, El\'ements de Math. \emph{27, 28, 30, 31} Hermann, Paris (1961--1965).\par
\noindent\hangindent=25bp\hangafter=1\hspace*{6bp}%
\makebox[19bp][r]{\hypertarget{bib-bourbaki-2}{}\hypertarget{bib-474-5}{}2.\enspace}\emph{Alg\`ebre}, El\'ements de Math. \emph{4, 6, 7, 11, 14, 23, 24}, Hermann, Paris (1947--59).\par
\par\vspace{2.15pt}\noindent \hypertarget{bib-474-6}{}Brieskorn, E.\par
\noindent\hangindent=25bp\hangafter=1\hspace*{6bp}%
\makebox[19bp][r]{\hypertarget{bib-474-7}{}1.\enspace}Beispiele zur Differentialtopologie von Singularit\"aten, \emph{Invent. Math. 2} (1966) 1--14.\par
\noindent\hangindent=25bp\hangafter=1\hspace*{6bp}%
\makebox[19bp][r]{\hypertarget{bib-brieskorn-2}{}2.\enspace}Rationale Singularit\"aten Komplexer Fl\"achen, \emph{Invent. Math. 4} (1968) 336--358.\par
\par\vspace{2.15pt}\noindent \hypertarget{bib-474-8}{}Burnside, W.\par
\noindent\hangindent=25bp\hangafter=1\hspace*{6bp}%
\makebox[19bp][r]{\hypertarget{bib-474-8-item-1}{}1.\enspace}\emph{Theory of Groups of Finite Order}, Cambridge Univ. Press, Cambridge (1911); reprinted by Dover, New York.\par
\par\vspace{2.15pt}\noindent \hypertarget{bib-474-9}{}Cartan, H. and Chevalley, C.\par
\noindent\hangindent=25bp\hangafter=1\hspace*{6bp}%
\makebox[19bp][r]{\hypertarget{bib-474-9-item-1}{}1.\enspace}\emph{G\'eom\'etrie Alg\'ebrique}, S\'eminaire Cartan--Chevalley, Secr\'etariat Math., Paris (1955/56).\par
\par\vspace{2.15pt}\noindent \hypertarget{bib-474-10}{}Cartan, H. and Eilenberg, S.\par
\noindent\hangindent=25bp\hangafter=1\hspace*{6bp}%
\makebox[19bp][r]{\hypertarget{bib-cartan-eilenberg-1}{}1.\enspace}\emph{Homological Algebra}, Princeton Univ. Press, Princeton (1956), xv + 390 pp.\par
\par\vspace{2.15pt}\noindent \hypertarget{bib-474-11}{}Cassels, J. W. S.\par
\noindent\hangindent=25bp\hangafter=1\hspace*{6bp}%
\makebox[19bp][r]{\hypertarget{bib-474-11-item-1}{}1.\enspace}Diophantine equations with special reference to elliptic curves, \emph{J. Lond. Math. Soc. 41} (1966), 193--291. Corr. \emph{42} (1967) 183.\par
\par\vspace{2.15pt}\noindent \hypertarget{bib-474-12}{}Cassels, J. W. S. and Fr\"ohlich, A., ed.\par
\noindent\hangindent=25bp\hangafter=1\hspace*{6bp}%
\makebox[19bp][r]{\hypertarget{bib-cassels-frohlich-1}{}1.\enspace}\emph{Algebraic Number Theory}, Thompson Book Co, Washington D.C. (1967).\par
\par\vspace{2.15pt}\noindent \hypertarget{bib-474-13}{}Castelnuovo, G.\par
\noindent\hangindent=25bp\hangafter=1\hspace*{6bp}%
\makebox[19bp][r]{\hypertarget{bib-474-13-item-1}{}1.\enspace}Sui multipli di una serie lineare di gruppi di punti appartenente ad una curva algebrica, \emph{Rend. Circ. Mat. Palermo 7} (1893). Also in \emph{Memorie Scelte}, pp.~95--113.\par
\par\vspace{2.15pt}\noindent \hypertarget{bib-474-14}{}Chevalley, C.\par
\noindent\hangindent=25bp\hangafter=1\hspace*{6bp}%
\makebox[19bp][r]{\hypertarget{bib-474-15}{}1.\enspace}Intersections of algebraic and algebroid varieties, \emph{Trans. Amer. Math. Soc. 57} (1945), 1--85.\par
\noindent\hangindent=25bp\hangafter=1\hspace*{6bp}%
\makebox[19bp][r]{\hypertarget{bib-chevalley-2}{}2.\enspace}\emph{Anneaux de Chow et Applications}, S\'eminaire Chevalley, Secr\'etariat Math., Paris (1958).\par
\par\vspace{2.15pt}\noindent \hypertarget{bib-474-16}{}Chow, W. L.\par
\noindent\hangindent=25bp\hangafter=1\hspace*{6bp}%
\makebox[19bp][r]{\hypertarget{bib-474-17}{}1.\enspace}On compact complex analytic varieties, \emph{Amer. J. of Math. 71} (1949), 893--914; errata \emph{72}, p.~624.\par
\noindent\hangindent=25bp\hangafter=1\hspace*{6bp}%
\makebox[19bp][r]{\hypertarget{bib-474-18}{}2.\enspace}On Picard varieties, \emph{Amer. J. Math. 74} (1952), 895--909.\par
\noindent\hangindent=25bp\hangafter=1\hspace*{6bp}%
\makebox[19bp][r]{\hypertarget{bib-chow-3}{}3.\enspace}The Jacobian variety of an algebraic curve, \emph{Amer. J. of Math. 76} (1954), 453--476.\par
\par\vspace{2.15pt}\noindent \hypertarget{bib-474-19}{}Chow, W. L. and Kodaira, K.\par
\noindent\hangindent=25bp\hangafter=1\hspace*{6bp}%
\makebox[19bp][r]{\hypertarget{bib-chow-kodaira-1}{}1.\enspace}On analytic surfaces with two independent meromorphic functions, \emph{Proc. Nat. Acad. Sci. USA 38} (1952), 319--325.\par
}
% Contents from body/page-475.tex.
% Original printed page 461.
\ChapterSheet[10.25]{475}{461}{\Goto{p-473}{Bibliography}}{%
\fontsize{9.25}{9.8}\selectfont%
\setlength{\parskip}{0pt}\setlength{\parindent}{0pt}%
\setlength{\emergencystretch}{1.7em}%
\par\vspace{2.15pt}\noindent \hypertarget{bib-475-0}{}Clemens, C. H. and Griffiths, P. A.\par
\noindent\hangindent=25bp\hangafter=1\hspace*{6bp}%
\makebox[19bp][r]{\hypertarget{bib-clemens-griffiths-1}{}1.\enspace}The intermediate Jacobian of the cubic threefold, \emph{Annals of Math. 95} (1972), 281--356.\par
\par\vspace{2.15pt}\noindent \hypertarget{bib-475-1}{}Deligne, P.\par
\noindent\hangindent=25bp\hangafter=1\hspace*{6bp}%
\makebox[19bp][r]{\hypertarget{bib-475-2}{}1.\enspace}Formes modulaires et repr\'esentations $l$-adiques, \emph{S\'eminaire Bourbaki 355}, Lecture Notes in Math. \emph{179}, Springer-Verlag, Heidelberg (1971), 139--172.\par
\noindent\hangindent=25bp\hangafter=1\hspace*{6bp}%
\makebox[19bp][r]{\hypertarget{bib-475-3}{}2.\enspace}La conjecture de Weil pour les surfaces K3, \emph{Invent. Math. 15} (1972), 206--226.\par
\noindent\hangindent=25bp\hangafter=1\hspace*{6bp}%
\makebox[19bp][r]{\hypertarget{bib-475-4}{}3.\enspace}La conjecture de Weil, I, \emph{Publ. Math. IHES 43} (1974), 273--307.\par
\noindent\hangindent=25bp\hangafter=1\hspace*{6bp}%
\makebox[19bp][r]{\hypertarget{bib-475-5}{}\hypertarget{bib-deligne-4}{}4.\enspace}Th\'eor\`emes de Lefschetz et crit\`eres de d\'eg\'en\'erescence de suites spectrales, \emph{Publ. Math. IHES 35} (1968) 107--126.\par
\noindent\hangindent=25bp\hangafter=1\hspace*{6bp}%
\makebox[19bp][r]{\hypertarget{bib-deligne-5}{}5.\enspace}Les intersections compl\`etes de niveau de Hodge un, \emph{Invent. Math. 15} (1972) 237--250.\par
\par\vspace{2.15pt}\noindent \hypertarget{bib-475-6}{}Deligne, P. and Mumford, D.\par
\noindent\hangindent=25bp\hangafter=1\hspace*{6bp}%
\makebox[19bp][r]{\hypertarget{bib-475-6-item-1}{}1.\enspace}The irreducibility of the space of curves of given genus, \emph{Publ. Math. IHES 36} (1969), 75--110.\par
\par\vspace{2.15pt}\noindent \hypertarget{bib-475-7}{}Demazure, M.\par
\noindent\hangindent=25bp\hangafter=1\hspace*{6bp}%
\makebox[19bp][r]{\hypertarget{bib-demazure-1}{}1.\enspace}Motifs des vari\'et\'es alg\'ebriques, \emph{S\'eminaire Bourbaki 365}, Lecture Notes in Math. \emph{180}, Springer-Verlag, Heidelberg (1971), 19--38.\par
\par\vspace{2.15pt}\noindent \hypertarget{bib-475-8}{}Deuring, M.\par
\noindent\hangindent=25bp\hangafter=1\hspace*{6bp}%
\makebox[19bp][r]{\hypertarget{bib-475-9}{}1.\enspace}Die Typen der Multiplikatorenringe elliptischer Funktionenk\"orper, \emph{Abh. Math. Sem. Univ. Hamburg 14} (1941), 197--272.\par
\noindent\hangindent=25bp\hangafter=1\hspace*{6bp}%
\makebox[19bp][r]{\hypertarget{bib-deuring-2}{}2.\enspace}\emph{Die Klassenk\"orper der Komplexen Multiplikation}, Enz. der Math. Wiss., 2nd ed., $\mathrm I_2$, Heft 10, II, \S23 (1958) 60 pp.\par
\par\vspace{2.15pt}\noindent \hypertarget{bib-475-10}{}Dieudonn\'e, J.\par
\noindent\hangindent=25bp\hangafter=1\hspace*{6bp}%
\makebox[19bp][r]{1.\enspace}\emph{Cours de G\'eom\'etrie Alg\'ebrique, I. Aper\c{c}u Historique sur le D\'eveloppement de la G\'eom\'etrie Alg\'ebrique}, Presses Univ. France, Collection Sup. (1974), 234 pp.\par
\par\vspace{2.15pt}\noindent \hypertarget{bib-475-11}{}Dwork, B.\par
\noindent\hangindent=25bp\hangafter=1\hspace*{6bp}%
\makebox[19bp][r]{\hypertarget{bib-dwork-1}{}1.\enspace}On the rationality of the zeta function of an algebraic variety, \emph{Amer. J. Math. 82} (1960), 631--648.\par
\par\vspace{2.15pt}\noindent \hypertarget{bib-475-12}{}Freyd, P.\par
\noindent\hangindent=25bp\hangafter=1\hspace*{6bp}%
\makebox[19bp][r]{\hypertarget{bib-freyd-1}{}1.\enspace}\emph{Abelian Categories, an Introduction to the Theory of Functors}, Harper \& Row, New York (1964), 164 pp.\par
\par\vspace{2.15pt}\noindent \hypertarget{bib-475-13}{}Fulton, W.\par
\noindent\hangindent=25bp\hangafter=1\hspace*{6bp}%
\makebox[19bp][r]{\hypertarget{bib-475-14}{}\hypertarget{bib-fulton-1}{}1.\enspace}\emph{Algebraic Curves}, W. A. Benjamin, New York (1969), xii + 226 pp.\par
\noindent\hangindent=25bp\hangafter=1\hspace*{6bp}%
\makebox[19bp][r]{\hypertarget{bib-fulton-2}{}2.\enspace}Riemann--Roch for singular varieties, in \emph{Algebraic Geometry, Arcata 1974}, Amer. Math. Soc. Proc. Symp. Pure Math \emph{29} (1975), 449--457.\par
\par\vspace{2.15pt}\noindent \hypertarget{bib-475-15}{}Godement, R.\par
\noindent\hangindent=25bp\hangafter=1\hspace*{6bp}%
\makebox[19bp][r]{\hypertarget{bib-godement-1}{}1.\enspace}\emph{Topologie Alg\'ebrique et Th\'eorie des Faisceaux}, Hermann, Paris (1958).\par
\par\vspace{2.15pt}\noindent \hypertarget{bib-475-16}{}Grauert, H.\par
\noindent\hangindent=25bp\hangafter=1\hspace*{6bp}%
\makebox[19bp][r]{\hypertarget{bib-475-17}{}\hypertarget{bib-grauert-1}{}1.\enspace}Ein Theorem der analytischen Garbentheorie und die Modulr\"aume komplexer Strukturen, \emph{Pub. Math. IHES 5} (1960), 233--292.\par
\noindent\hangindent=25bp\hangafter=1\hspace*{6bp}%
\makebox[19bp][r]{\hypertarget{bib-grauert-2}{}2.\enspace}\"Uber Modifikationen und exzeptionelle analytische Mengen, \emph{Math. Ann. 146} (1962), 331--368.\par
\par\vspace{2.15pt}\noindent \hypertarget{bib-475-18}{}Grauert, H. and Remmert, R.\par
\noindent\hangindent=25bp\hangafter=1\hspace*{6bp}%
\makebox[19bp][r]{\hypertarget{bib-grauert-remmert-1}{}1.\enspace}Komplex R\"aume, \emph{Math. Ann. 136} (1958), 245--318.\par
\par\vspace{2.15pt}\noindent \hypertarget{bib-475-19}{}Grauert, H. and Riemenschneider, O.\par
\noindent\hangindent=25bp\hangafter=1\hspace*{6bp}%
\makebox[19bp][r]{\hypertarget{bib-grauert-riemenschneider-1}{}1.\enspace}Verschwindungss\"atze f\"ur analytische Kohomologiegruppen auf komplexen R\"aumen, \emph{Inv. Math. 11} (1970), 263--292.\par
}
% Contents from body/page-476.tex.
% Original printed page 462.
\ChapterSheet[10.25]{476}{462}{\Goto{p-473}{Bibliography}}{%
\fontsize{9.25}{9.8}\selectfont%
\setlength{\parskip}{0pt}\setlength{\parindent}{0pt}%
\setlength{\emergencystretch}{1.7em}%
\par\vspace{2.15pt}\noindent \hypertarget{bib-476-0}{}Grothendieck, A.\par
\noindent\hangindent=25bp\hangafter=1\hspace*{6bp}%
\makebox[19bp][r]{\hypertarget{bib-476-1}{}\hypertarget{bib-grothendieck-1}{}1.\enspace}Sur quelques points d'alg\`ebre homologique, \emph{T\^ohoku Math. J. 9} (1957), 119--221.\par
\noindent\hangindent=25bp\hangafter=1\hspace*{6bp}%
\makebox[19bp][r]{\hypertarget{bib-476-2}{}2.\enspace}Sur une note de Mattuck-Tate, \emph{J. Reine u. Angew. Math. 200} (1958), 208--215.\par
\noindent\hangindent=25bp\hangafter=1\hspace*{6bp}%
\makebox[19bp][r]{\hypertarget{bib-476-3}{}3.\enspace}La th\'eorie des classes de Chern, \emph{Bull. Soc. Math. de France 86} (1958), 137--154.\par
\noindent\hangindent=25bp\hangafter=1\hspace*{6bp}%
\makebox[19bp][r]{\hypertarget{bib-476-4}{}4.\enspace}The cohomology theory of abstract algebraic varieties, \emph{Proc. Int. Cong. Math., Edinburgh} (1958), 103--118.\par
\noindent\hangindent=25bp\hangafter=1\hspace*{6bp}%
\makebox[19bp][r]{\hypertarget{bib-476-5}{}5.\enspace}\emph{Fondements de la G\'eom\'etrie Alg\'ebrique}, S\'eminaire Bourbaki 1957--62, Secr\'etariat Math., Paris (1962).\par
\noindent\hangindent=25bp\hangafter=1\hspace*{6bp}%
\makebox[19bp][r]{\hypertarget{bib-476-6}{}6.\enspace}Formule de Lefschetz et rationalit\'e des fonctions $L$, \emph{S\'eminaire Bourbaki 279} (1965).\par
\noindent\hangindent=25bp\hangafter=1\hspace*{6bp}%
\makebox[19bp][r]{\hypertarget{bib-476-7}{}\hypertarget{bib-grothendieck-7}{}7.\enspace}\emph{Local Cohomology} (notes by R. Hartshorne), Lecture Notes in Math. \emph{41}, Springer-Verlag, Heidelberg (1967), 106 pp.\par
\noindent\hangindent=25bp\hangafter=1\hspace*{6bp}%
\makebox[19bp][r]{\hypertarget{bib-grothendieck-8}{}8.\enspace}Crystals and the De Rham cohomology of schemes (notes by I. Coates and O. Jussila), in \emph{Dix Expos\'es sur la Cohomologie des Sch\'emas}, North-Holland, Amsterdam (1968), 306--358.\par
\noindent\hangindent=25bp\hangafter=1\hspace*{6bp}%
\makebox[19bp][r]{\hypertarget{bib-476-8}{}9.\enspace}Standard conjectures on algebraic cycles, in \emph{Algebraic Geometry, Bombay 1968}, Oxford University Press, Oxford (1969), 193--199.\par
\par\vspace{2.15pt}\noindent Grothendieck, A. and Dieudonn\'e, J.\par
\noindent\hspace*{12bp}\hypertarget{bib-EGA}{} \emph{El\'ements de G\'eom\'etrie Alg\'ebrique}.\par
\noindent\hangindent=52bp\hangafter=1\hspace*{10bp}%
\makebox[42bp][l]{\hypertarget{bib-EGA-1}{}EGA I.}Le langage des sch\'emas, \emph{Publ. Math. IHES 4} (1960).\par
\noindent\hangindent=52bp\hangafter=1\hspace*{10bp}%
\makebox[42bp][l]{\hypertarget{bib-EGA-2}{}EGA II.}Etude globale \'el\'ementaire de quelques classes de morphismes, \emph{Ibid. 8} (1961).\par
\noindent\hangindent=52bp\hangafter=1\hspace*{10bp}%
\makebox[42bp][l]{\hypertarget{bib-EGA-3}{}\hypertarget{bib-ega-iii}{}EGA III.}Etude cohomologique des faisceaux coh\'erents, \emph{Ibid. 11} (1961), and \emph{17} (1963).\par
\noindent\hangindent=52bp\hangafter=1\hspace*{10bp}%
\makebox[42bp][l]{\hypertarget{bib-EGA-4}{}EGA IV.}Etude locale des sch\'emas et des morphismes de sch\'emas, \emph{Ibid. 20} (1964), \emph{24} (1965), \emph{28} (1966), \emph{32} (1967).\par
\noindent\hangindent=52bp\hangafter=1\hspace*{10bp}%
\makebox[42bp][l]{EGA I.}\emph{El\'ements de G\'eom\'etrie Alg\'ebrique, I}, Grundlehren \emph{166}, Springer-Verlag, Heidelberg (new ed., 1971), ix + 466 pp.\par
\par\vspace{2.15pt}\noindent Grothendieck, A. et al.\par
\noindent\hspace*{12bp}\hypertarget{bib-SGA}{} \emph{S\'eminaire de G\'eom\'etrie Alg\'ebrique}.\par
\noindent\hangindent=52bp\hangafter=1\hspace*{10bp}%
\makebox[42bp][l]{\hypertarget{bib-SGA-1}{}SGA 1.}\emph{Rev\^etements \'etales et Groupe Fondemental}, Lecture Notes in Math. \emph{224}, Springer-Verlag, Heidelberg (1971).\par
\noindent\hangindent=52bp\hangafter=1\hspace*{10bp}%
\makebox[42bp][l]{\hypertarget{bib-SGA-2}{}SGA 2.}\emph{Cohomologie Locale des Faisceaux Coh\'erents et Th\'eor\`emes de Lefschetz Locaux et Globaux}, North-Holland, Amsterdam (1968).\par
\noindent\hangindent=52bp\hangafter=1\hspace*{10bp}%
\makebox[42bp][l]{\hypertarget{bib-SGA-3}{}SGA 3.}(with Demazure, M.) \emph{Sch\'emas en Groupes I, II, III}, Lecture Notes in Math. \emph{151, 152, 153}, Springer-Verlag, Heidelberg (1970).\par
\noindent\hangindent=52bp\hangafter=1\hspace*{10bp}%
\makebox[42bp][l]{\hypertarget{bib-SGA-4}{}SGA 4.}(with Artin, M. and Verdier, J. L.) \emph{Th\'eorie des Topos et Cohomologie \'Etale des Sch\'emas}, Lecture Notes in Math. \emph{269, 270, 305}, Springer-Verlag, Heidelberg (1972--1973).\par
\noindent\hangindent=52bp\hangafter=1\hspace*{10bp}%
\makebox[42bp][l]{\hypertarget{bib-SGA-4half}{}SGA $4\frac12$.}(by Deligne, P., with Boutot, J. F., Illusie, L., and Verdier, J. L.) \emph{Cohomologie Etale}, Lecture Notes in Math. \emph{569}, Springer-Verlag, Heidelberg (1977).\par
\noindent\hangindent=52bp\hangafter=1\hspace*{10bp}%
\makebox[42bp][l]{\hypertarget{bib-SGA-5}{}SGA 5.}\emph{Cohomologie $l$-adique et fonctions $L$} (unpublished).\par
\noindent\hangindent=52bp\hangafter=1\hspace*{10bp}%
\makebox[42bp][l]{\hypertarget{bib-SGA-6}{}SGA 6.}(with Berthelot, P. and Illusie, L.) \emph{Th\'eorie des Intersections et Th\'eor\`eme de Riemann--Roch}, Lecture Notes in Math. \emph{225}, Springer-Verlag, Heidelberg (1971).\par
\noindent\hangindent=52bp\hangafter=1\hspace*{10bp}%
\makebox[42bp][l]{\hypertarget{bib-SGA-7}{}SGA 7.}(with Raynaud, M. and Rim, D. S.) \emph{Groupes de Monodromie en G\'eom\'etrie Alg\'ebrique}, Lecture Notes in Math. \emph{288}, Springer-Verlag, Heidelberg (1972). Part II (by Deligne, P. and Katz, N.) \emph{340} (1973).\par
\par\vspace{2.15pt}\noindent \hypertarget{bib-476-19}{}Gunning, R. C.\par
\noindent\hangindent=25bp\hangafter=1\hspace*{6bp}%
\makebox[19bp][r]{\hypertarget{bib-gunning-1}{}1.\enspace}\emph{Lectures on Riemann Surfaces}, Princeton Math. Notes, Princeton U. Press, Princeton (1966), 256 pp.\par
}
% Contents from body/page-477.tex.
% Original printed page 463.
\ChapterSheet[10.25]{477}{463}{\Goto{p-473}{Bibliography}}{%
\fontsize{9.25}{9.8}\selectfont%
\setlength{\parskip}{0pt}\setlength{\parindent}{0pt}%
\setlength{\emergencystretch}{1.7em}%
\par\vspace{2.15pt}\noindent \hypertarget{bib-477-0}{}Gunning, R. C. and Rossi, H.\par
\noindent\hangindent=25bp\hangafter=1\hspace*{6bp}%
\makebox[19bp][r]{\hypertarget{bib-gunning-rossi-1}{}1.\enspace}\emph{Analytic Functions of Several Complex Variables}, Prentice-Hall (1965), xii + 317 pp.\par
\par\vspace{2.15pt}\noindent \hypertarget{bib-477-1}{}Halphen, G.\par
\noindent\hangindent=25bp\hangafter=1\hspace*{6bp}%
\makebox[19bp][r]{\hypertarget{bib-477-1-item-1}{}1.\enspace}M\'emoire sur la classification des courbes gauches alg\'ebriques, \emph{J. \'Ec. Polyt. 52} (1882), 1--200.\par
\par\vspace{2.15pt}\noindent \hypertarget{bib-477-2}{}Hartshorne, R.\par
\noindent\hangindent=25bp\hangafter=1\hspace*{6bp}%
\makebox[19bp][r]{\hypertarget{bib-477-3}{}\hypertarget{legacy-b-h1}{}1.\enspace}Complete intersections and connectedness, \emph{Amer. J. of Math. 84} (1962), 497--508.\par
\noindent\hangindent=25bp\hangafter=1\hspace*{6bp}%
\makebox[19bp][r]{\hypertarget{bib-hartshorne-2}{}2.\enspace}\emph{Residues and Duality}, Lecture Notes in Math. \emph{20}, Springer-Verlag, Heidelberg (1966).\par
\noindent\hangindent=25bp\hangafter=1\hspace*{6bp}%
\makebox[19bp][r]{\hypertarget{bib-477-4}{}\hypertarget{bib-hartshorne-3}{}3.\enspace}Cohomological dimension of algebraic varieties, \emph{Annals of Math. 88} (1968), 403--450.\par
\noindent\hangindent=25bp\hangafter=1\hspace*{6bp}%
\makebox[19bp][r]{\hypertarget{bib-477-5}{}4.\enspace}Curves with high self-intersection on algebraic surfaces, \emph{Publ. Math. IHES 36} (1969), 111--125.\par
\noindent\hangindent=25bp\hangafter=1\hspace*{6bp}%
\makebox[19bp][r]{\hypertarget{bib-477-6}{}\hypertarget{bib-hartshorne-5}{}\hypertarget{legacy-b-h5}{}5.\enspace}\emph{Ample Subvarieties of Algebraic Varieties}, Lecture Notes in Math. \emph{156}, Springer-Verlag, Heidelberg (1970), xiii + 256 pp.\par
\noindent\hangindent=25bp\hangafter=1\hspace*{6bp}%
\makebox[19bp][r]{\hypertarget{bib-477-7}{}\hypertarget{bib-hartshorne-6}{}6.\enspace}Equivalence relations on algebraic cycles and subvarieties of small codimension, in \emph{Algebraic Geometry, Arcata 1974}, Amer. Math. Soc. Proc. Symp. Pure Math. \emph{29} (1975), 129--164.\par
\noindent\hangindent=25bp\hangafter=1\hspace*{6bp}%
\makebox[19bp][r]{\hypertarget{bib-477-8}{}7.\enspace}On the De Rham cohomology of algebraic varieties, \emph{Publ. Math. IHES 45} (1976), 5--99.\par
\par\vspace{2.15pt}\noindent \hypertarget{bib-477-9}{}Hartshorne, R., ed.\par
\noindent\hangindent=25bp\hangafter=1\hspace*{6bp}%
\makebox[19bp][r]{1.\enspace}\emph{Algebraic Geometry, Arcata 1974}, Amer. Math. Soc. Proc. Symp. Pure Math. \emph{29} (1975).\par
\par\vspace{2.15pt}\noindent \hypertarget{bib-477-10}{}Hilbert, D. and Cohn-Vossen, S.\par
\noindent\hangindent=25bp\hangafter=1\hspace*{6bp}%
\makebox[19bp][r]{\hypertarget{bib-477-10-item-1}{}1.\enspace}\emph{Geometry and the Imagination}, Chelsea Pub. Co., New York (1952), ix + 357 pp. (translated from German \emph{Anschauliche Geometrie} (1932)).\par
\par\vspace{2.15pt}\noindent \hypertarget{bib-477-11}{}Hilton, P. J. and Stammbach, U.\par
\noindent\hangindent=25bp\hangafter=1\hspace*{6bp}%
\makebox[19bp][r]{\hypertarget{bib-hilton-stammbach-1}{}1.\enspace}\emph{A Course in Homological Algebra}, Graduate Texts in Mathematics \emph{4}, Springer-Verlag, Heidelberg (1970), ix + 338 pp.\par
\par\vspace{2.15pt}\noindent \hypertarget{bib-477-12}{}Hironaka, H.\par
\noindent\hangindent=25bp\hangafter=1\hspace*{6bp}%
\makebox[19bp][r]{\hypertarget{bib-477-12-item-1}{}1.\enspace}A note on algebraic geometry over ground rings. The invariance of Hilbert characteristic functions under the specialization process, \emph{Ill. J. Math. 2} (1958), 355--366.\par
\noindent\hangindent=25bp\hangafter=1\hspace*{6bp}%
\makebox[19bp][r]{\hypertarget{bib-477-13}{}\hypertarget{bib-477-14}{}2.\enspace}On the theory of birational blowing-up, Thesis, Harvard (1960) (unpublished).\par
\noindent\hangindent=25bp\hangafter=1\hspace*{6bp}%
\makebox[19bp][r]{\hypertarget{bib-477-15}{}3.\enspace}On resolution of singularities (characteristic zero), \emph{Proc. Int. Cong. Math.} (1962), 507--521.\par
\noindent\hangindent=25bp\hangafter=1\hspace*{6bp}%
\makebox[19bp][r]{\hypertarget{bib-477-16}{}\hypertarget{bib-hironaka-4}{}4.\enspace}Resolution of singularities of an algebraic variety over a field of characteristic zero, \emph{Annals of Math. 79} (1964). I: 109--203; II: 205--326.\par
\noindent\hangindent=25bp\hangafter=1\hspace*{6bp}%
\makebox[19bp][r]{\hypertarget{bib-hironaka-5}{}5.\enspace}Triangulations of algebraic sets, in \emph{Algebraic Geometry, Arcata 1974}, Amer. Math. Soc. Proc. Symp. Pure Math. \emph{29} (1975), 165--184.\par
\par\vspace{2.15pt}\noindent \hypertarget{bib-477-17}{}Hironaka, H. and Matsumura, H.\par
\noindent\hangindent=25bp\hangafter=1\hspace*{6bp}%
\makebox[19bp][r]{\hypertarget{bib-477-17-item-1}{}1.\enspace}Formal functions and formal embeddings, \emph{J. of Math. Soc. of Japan 20} (1968), 52--82.\par
\par\vspace{2.15pt}\noindent \hypertarget{bib-477-18}{}Hirzebruch, F.\par
\noindent\hangindent=25bp\hangafter=1\hspace*{6bp}%
\makebox[19bp][r]{\hypertarget{bib-hirzebruch-1}{}1.\enspace}\emph{Topological Methods in Algebraic Geometry}, Grundlehren \emph{131}, Springer-Verlag, Heidelberg (3rd ed., 1966), ix + 232 pp.\par
\par\vspace{2.15pt}\noindent \hypertarget{bib-477-19}{}Hirzebruch, F. and Mayer, K. H.\par
\noindent\hangindent=25bp\hangafter=1\hspace*{6bp}%
\makebox[19bp][r]{\hypertarget{bib-477-19-item-1}{}1.\enspace}\emph{O(n)-Mannigfaltigkeiten, Exotische Sph\"aren und Singularit\"aten}, Lecture Notes in Math. \emph{57}, Springer-Verlag, Heidelberg (1968).\par
}
% Contents from body/page-478.tex.
% Original printed page 464.
\ChapterSheet[10.25]{478}{464}{\Goto{p-473}{Bibliography}}{%
\fontsize{9.25}{9.8}\selectfont%
\setlength{\parskip}{0pt}\setlength{\parindent}{0pt}%
\setlength{\emergencystretch}{1.7em}%
\par\vspace{2.15pt}\noindent \hypertarget{bib-478-0}{}Hodge, W. V. D.\par
\noindent\hangindent=25bp\hangafter=1\hspace*{6bp}%
\makebox[19bp][r]{\hypertarget{bib-hodge-1}{}1.\enspace}\emph{The Theory and Applications of Harmonic Integrals}, Cambridge Univ. Press, Cambridge (2nd ed., 1952), 282 pp.\par
\par\vspace{2.15pt}\noindent \hypertarget{bib-478-1}{}Horrocks, G. and Mumford, D.\par
\noindent\hangindent=25bp\hangafter=1\hspace*{6bp}%
\makebox[19bp][r]{\hypertarget{bib-horrocks-mumford-1}{}1.\enspace}A rank 2 vector bundle on $\mathbf P^4$ with 15,000 symmetries, \emph{Topology 12} (1973), 63--81.\par
\par\vspace{2.15pt}\noindent \hypertarget{bib-478-2}{}Hurwitz, A.\par
\noindent\hangindent=25bp\hangafter=1\hspace*{6bp}%
\makebox[19bp][r]{\hypertarget{bib-hurwitz-1}{}1.\enspace}\"Uber algebraische Gebilde mit eindeutigen Transformationen in sich, \emph{Math. Ann. 41} (1893), 403--442.\par
\par\vspace{2.15pt}\noindent \hypertarget{bib-478-3}{}Hurwitz, A. and Courant, R.\par
\noindent\hangindent=25bp\hangafter=1\hspace*{6bp}%
\makebox[19bp][r]{\hypertarget{bib-hurwitz-courant-1}{}1.\enspace}\emph{Allgemeine Funktionentheorie und elliptische Funktionen; Geometrische Funktionentheorie}, Grundlehren \emph{3}, Springer-Verlag, Heidelberg (1922), xi + 399 pp.\par
\par\vspace{2.15pt}\noindent \hypertarget{bib-478-4}{}Igusa, J.-I.\par
\noindent\hangindent=25bp\hangafter=1\hspace*{6bp}%
\makebox[19bp][r]{\hypertarget{bib-478-4-item-1}{}1.\enspace}Arithmetic genera of normal varieties in an algebraic family, \emph{Proc. Nat. Acad. Sci. USA 41} (1955) 34--37.\par
\noindent\hangindent=25bp\hangafter=1\hspace*{6bp}%
\makebox[19bp][r]{\hypertarget{bib-478-5}{}2.\enspace}Class number of a definite quaternion with prime discriminant, \emph{Proc. Nat. Acad. Sci. USA 44} (1958) 312--314.\par
\par\vspace{2.15pt}\noindent \hypertarget{bib-478-6}{}Iskovskih, V. A. and Manin, Ju. I.\par
\noindent\hangindent=25bp\hangafter=1\hspace*{6bp}%
\makebox[19bp][r]{\hypertarget{bib-iskovskih-manin-1}{}1.\enspace}Three-dimensional quartics and counterexamples to the L\"uroth problem, \emph{Math. USSR---Sbornik 15} (1971) 141--166.\par
\par\vspace{2.15pt}\noindent \hypertarget{bib-478-7}{}Jouanolou, J. P.\par
\noindent\hangindent=25bp\hangafter=1\hspace*{6bp}%
\makebox[19bp][r]{\hypertarget{bib-jouanolou-1}{}1.\enspace}Riemann--Roch sans d\'enominateurs, \emph{Invent. Math. 11} (1970), 15--26.\par
\par\vspace{2.15pt}\noindent \hypertarget{bib-478-8}{}Kleiman, S. L.\par
\noindent\hangindent=25bp\hangafter=1\hspace*{6bp}%
\makebox[19bp][r]{\hypertarget{bib-478-9}{}1.\enspace}Toward a numerical theory of ampleness, \emph{Annals of Math. 84} (1966), 293--344.\par
\noindent\hangindent=25bp\hangafter=1\hspace*{6bp}%
\makebox[19bp][r]{\hypertarget{bib-478-10}{}2.\enspace}Algebraic cycles and the Weil conjectures, in \emph{Dix Expos\'es sur la Cohomologie des Sch\'emas}, North-Holland, Amsterdam (1968), 359--386.\par
\noindent\hangindent=25bp\hangafter=1\hspace*{6bp}%
\makebox[19bp][r]{\hypertarget{bib-kleiman-3}{}3.\enspace}The transversality of a general translate, \emph{Compos. Math. 28} (1974), 287--297.\par
\par\vspace{2.15pt}\noindent \hypertarget{bib-478-11}{}Kleiman, S. L. and Laksov, D.\par
\noindent\hangindent=25bp\hangafter=1\hspace*{6bp}%
\makebox[19bp][r]{\hypertarget{bib-478-11-item-1}{}1.\enspace}Another proof of the existence of special divisors, \emph{Acta Math. 132} (1974), 163--176.\par
\par\vspace{2.15pt}\noindent \hypertarget{bib-478-12}{}Klein, F.\par
\noindent\hangindent=25bp\hangafter=1\hspace*{6bp}%
\makebox[19bp][r]{\hypertarget{bib-478-13}{}1.\enspace}Ueber die Transformationen siebenter Ordnung der elliptischen Funktionen, \emph{Math. Ann. 14} (1879). Also in Klein, \emph{Ges. Math. Abh. 3} (1923), 90--136.\par
\noindent\hangindent=25bp\hangafter=1\hspace*{6bp}%
\makebox[19bp][r]{\hypertarget{bib-klein-2}{}2.\enspace}\emph{Lectures on the Icosahedron and the Solution of Equations of the Fifth Degree}, Kegan Paul, Trench, Tr\"ubner, London (1913). Dover reprint (1956).\par
\par\vspace{2.15pt}\noindent \hypertarget{bib-478-14}{}Knutson, D.\par
\noindent\hangindent=25bp\hangafter=1\hspace*{6bp}%
\makebox[19bp][r]{\hypertarget{bib-knutson-1}{}1.\enspace}\emph{Algebraic spaces}, Lecture Notes in Math. \emph{203}, Springer-Verlag, Heidelberg (1971) vi + 261 pp.\par
\par\vspace{2.15pt}\noindent \hypertarget{bib-478-15}{}Kodaira, K.\par
\noindent\hangindent=25bp\hangafter=1\hspace*{6bp}%
\makebox[19bp][r]{\hypertarget{bib-kodaira-1}{}1.\enspace}On a differential-geometric method in the theory of analytic stacks, \emph{Proc. Nat. Acad. Sci. USA 39} (1953), 1268--1273.\par
\noindent\hangindent=25bp\hangafter=1\hspace*{6bp}%
\makebox[19bp][r]{\hypertarget{bib-478-16}{}2.\enspace}On K\"ahler varieties of restricted type. (An intrinsic characterization of algebraic varieties.), \emph{Annals of Math. 60} (1954), 28--48.\par
\par\vspace{2.15pt}\noindent \hypertarget{bib-478-17}{}Kodaira, K. and Spencer, D. C.\par
\noindent\hangindent=25bp\hangafter=1\hspace*{6bp}%
\makebox[19bp][r]{\hypertarget{bib-478-17-item-1}{}1.\enspace}On arithmetic genera of algebraic varieties, \emph{Proc. Nat. Acad. Sci. USA 39} (1953), 641--649.\par
\par\vspace{2.15pt}\noindent \hypertarget{bib-478-18}{}Kunz, E.\par
\noindent\hangindent=25bp\hangafter=1\hspace*{6bp}%
\makebox[19bp][r]{\hypertarget{bib-478-18-item-1}{}1.\enspace}Holomorphe Differentialformen auf algebraischen Variet\"aten mit Singularit\"aten, I, \emph{Manus. Math. 15} (1975) 91--108.\par
}
% Contents from body/page-479.tex.
% Original printed page 465.
\ChapterSheet[10.25]{479}{465}{\Goto{p-473}{Bibliography}}{%
\fontsize{9.25}{9.8}\selectfont%
\setlength{\parskip}{0pt}\setlength{\parindent}{0pt}%
\setlength{\emergencystretch}{1.7em}%
\par\vspace{2.15pt}\noindent \hypertarget{bib-479-0}{}Lang, S.\par
\noindent\hangindent=25bp\hangafter=1\hspace*{6bp}%
\makebox[19bp][r]{\hypertarget{bib-479-1}{}\hypertarget{bib-lang-1}{}\hypertarget{legacy-b-lang}{}1.\enspace}\emph{Abelian Varieties}, Interscience Pub., New York (1959), xii + 256 pp.\par
\noindent\hangindent=25bp\hangafter=1\hspace*{6bp}%
\makebox[19bp][r]{\hypertarget{bib-lang-2}{}2.\enspace}\emph{Algebra}, Addison-Wesley (1971), xvii + 526 pp.\par
\par\vspace{2.15pt}\noindent \hypertarget{bib-479-2}{}Lang, S. and N\'eron, A.\par
\noindent\hangindent=25bp\hangafter=1\hspace*{6bp}%
\makebox[19bp][r]{\hypertarget{bib-lang-neron-1}{}1.\enspace}Rational points of abelian varieties over function fields, \emph{Amer. J. Math. 81} (1959), 95--118.\par
\par\vspace{2.15pt}\noindent \hypertarget{bib-479-3}{}Lang, S. and Trotter, H.\par
\noindent\hangindent=25bp\hangafter=1\hspace*{6bp}%
\makebox[19bp][r]{\hypertarget{bib-479-3-item-1}{}1.\enspace}\emph{Frobenius Distributions in} $\mathrm{GL}_2$\emph{-Extensions}, Lecture Notes in Math. \emph{504}, Springer-Verlag, Heidelberg (1976).\par
\par\vspace{2.15pt}\noindent \hypertarget{bib-479-4}{}Lang, S. and Weil, A.\par
\noindent\hangindent=25bp\hangafter=1\hspace*{6bp}%
\makebox[19bp][r]{\hypertarget{bib-lang-weil-1}{}1.\enspace}Number of points of varieties in finite fields, \emph{Amer. J. Math. 76} (1954), 819--827.\par
\par\vspace{2.15pt}\noindent \hypertarget{bib-479-5}{}Lascu, A. T., Mumford, D., and Scott, D. B.\par
\noindent\hangindent=25bp\hangafter=1\hspace*{6bp}%
\makebox[19bp][r]{\hypertarget{bib-lascu-mumford-scott-1}{}1.\enspace}The self-intersection formula and the \textquotedblleft formule-clef,\textquotedblright\ \emph{Math. Proc. Camb. Phil. Soc. 78} (1975), 117--123.\par
\par\vspace{2.15pt}\noindent \hypertarget{bib-479-6}{}Lichtenbaum, S. and Schlessinger, M.\par
\noindent\hangindent=25bp\hangafter=1\hspace*{6bp}%
\makebox[19bp][r]{\hypertarget{bib-479-6-item-1}{}1.\enspace}The cotangent complex of a morphism, \emph{Trans. Amer. Math. Soc. 128} (1967), 41--70.\par
\par\vspace{2.15pt}\noindent \hypertarget{bib-479-7}{}Lipman, J.\par
\noindent\hangindent=25bp\hangafter=1\hspace*{6bp}%
\makebox[19bp][r]{\hypertarget{bib-479-7-item-1}{}1.\enspace}Introduction to resolution of singularities, in \emph{Algebraic Geometry, Arcata 1974}, Amer. Math. Soc. Proc. Symp. Pure Math. \emph{29} (1975), 187--230.\par
\noindent\hangindent=25bp\hangafter=1\hspace*{6bp}%
\makebox[19bp][r]{\hypertarget{bib-479-8}{}\hypertarget{bib-479-9}{}2.\enspace}Rational singularities with applications to algebraic surfaces and unique factorization, \emph{Publ. Math. IHES 36} (1969) 195--279\par
\noindent\hangindent=25bp\hangafter=1\hspace*{6bp}%
\makebox[19bp][r]{\hypertarget{bib-lipman-3}{}3.\enspace}Unique factorization in complete local rings, in \emph{Algebraic Geometry, Arcata 1974}, Amer. Math. Soc. Proc. Symp. Pure Math. \emph{29} (1975) 531--546.\par
\par\vspace{2.15pt}\noindent \hypertarget{bib-479-10}{}Lubkin, S.\par
\noindent\hangindent=25bp\hangafter=1\hspace*{6bp}%
\makebox[19bp][r]{\hypertarget{bib-lubkin-1}{}1.\enspace}A $p$-adic proof of Weil's conjectures, \emph{Annals of Math. 87} (1968), 105--194, and \emph{87} (1968), 195--255.\par
\par\vspace{2.15pt}\noindent \hypertarget{bib-479-11}{}Macbeath, A. M.\par
\noindent\hangindent=25bp\hangafter=1\hspace*{6bp}%
\makebox[19bp][r]{\hypertarget{bib-macbeath-1}{}1.\enspace}On a theorem of Hurwitz, \emph{Proc. Glasgow Math. Assoc. 5} (1961) 90--96.\par
\par\vspace{2.15pt}\noindent \hypertarget{bib-479-12}{}Manin, Yu. I.\par
\noindent\hangindent=25bp\hangafter=1\hspace*{6bp}%
\makebox[19bp][r]{\hypertarget{bib-manin-1}{}1.\enspace}\emph{Lectures on the $K$-functor in Algebraic Geometry}, Russian Mathematical Surveys \emph{24} (5) (1969), 1--89.\par
\noindent\hangindent=25bp\hangafter=1\hspace*{6bp}%
\makebox[19bp][r]{\hypertarget{bib-479-13}{}2.\enspace}Correspondences, motifs, and monoidal transformations, \emph{Math. USSR---Sbornik 6} (1968) 439--470.\par
\noindent\hangindent=25bp\hangafter=1\hspace*{6bp}%
\makebox[19bp][r]{\hypertarget{bib-479-14}{}3.\enspace}\emph{Cubic forms: Algebra, Geometry, Arithmetic}, North-Holland, Amsterdam (1974), vii + 292 pp.\par
\par\vspace{2.15pt}\noindent \hypertarget{bib-479-15}{}Maruyama, M.\par
\noindent\hangindent=25bp\hangafter=1\hspace*{6bp}%
\makebox[19bp][r]{\hypertarget{bib-479-15-item-1}{}1.\enspace}\emph{On Classification of Ruled Surfaces}, Kyoto Univ., Lectures in Math. \emph{3}, Kinokuniya, Tokyo (1970).\par
\par\vspace{2.15pt}\noindent \hypertarget{bib-479-16}{}Matsumura, H.\par
\noindent\hangindent=25bp\hangafter=1\hspace*{6bp}%
\makebox[19bp][r]{\hypertarget{bib-479-16-item-1}{}1.\enspace}Geometric structure of the cohomology rings in abstract algebraic geometry, \emph{Mem. Coll. Sci. Univ. Kyoto (A) 32} (1959), 33--84.\par
\noindent\hangindent=25bp\hangafter=1\hspace*{6bp}%
\makebox[19bp][r]{\hypertarget{bib-479-17}{}\hypertarget{bib-matsumura-2}{}\hypertarget{legacy-b-mat}{}2.\enspace}\emph{Commutative Algebra}, W. A. Benjamin Co., New York (1970), xii + 262 pp.\par
\par\vspace{2.15pt}\noindent \hypertarget{bib-479-18}{}Mattuck, A. and Tate, J.\par
\noindent\hangindent=25bp\hangafter=1\hspace*{6bp}%
\makebox[19bp][r]{\hypertarget{bib-mattuck-tate-1}{}1.\enspace}On the inequality of Castelnuovo-Severi, \emph{Abh. Math. Sem. Univ. Hamburg 22} (1958), 295--299.\par
\par\vspace{2.15pt}\noindent \hypertarget{bib-479-19}{}Moishezon, B. G.\par
\noindent\hangindent=25bp\hangafter=1\hspace*{6bp}%
\makebox[19bp][r]{\hypertarget{bib-moishezon-1}{}1.\enspace}A criterion for projectivity of complete algebraic abstract varieties, \emph{Amer. Math. Soc. Translations 63} (1967), 1--50.\par
}
% Contents from body/page-480.tex.
% Original printed page 466.
\ChapterSheet[10.25]{480}{466}{\Goto{p-473}{Bibliography}}{%
\fontsize{9.25}{9.8}\selectfont%
\setlength{\parskip}{0pt}\setlength{\parindent}{0pt}%
\setlength{\emergencystretch}{1.7em}%
\noindent\hangindent=25bp\hangafter=1\hspace*{6bp}%
\makebox[19bp][r]{\hypertarget{bib-480-0}{}2.\enspace}On $n$-dimensional compact varieties with $n$ algebraically independent meromorphic functions, \emph{Amer. Math. Soc. Translations 63} (1967), 51--177.\par
\par\vspace{2.15pt}\noindent \hypertarget{bib-480-1}{}Morrow, J. and Kodaira, K.\par
\noindent\hangindent=25bp\hangafter=1\hspace*{6bp}%
\makebox[19bp][r]{\hypertarget{bib-morrow-kodaira-1}{}1.\enspace}\emph{Complex Manifolds}, Holt, Rinehart \& Winston, New York (1971), vii + 192 pp.\par
\par\vspace{2.15pt}\noindent \hypertarget{bib-480-2}{}Mumford, D.\par
\noindent\hangindent=25bp\hangafter=1\hspace*{6bp}%
\makebox[19bp][r]{\hypertarget{bib-480-3}{}1.\enspace}\emph{Geometric Invariant Theory}, Ergebnisse, Springer-Verlag, Heidelberg (1965), vi + 146 pp.\par
\noindent\hangindent=25bp\hangafter=1\hspace*{6bp}%
\makebox[19bp][r]{\hypertarget{bib-480-4}{}\hypertarget{bib-mumford-2}{}2.\enspace}\emph{Lectures on Curves on an Algebraic Surface}, Annals of Math. Studies \emph{59}, Princeton U. Press, Princeton (1966).\par
\noindent\hangindent=25bp\hangafter=1\hspace*{6bp}%
\makebox[19bp][r]{\hypertarget{bib-mumford-3}{}\hypertarget{bib-480-5}{}3.\enspace}Pathologies, III, \emph{Amer. J. Math. 89} (1967), 94--104.\par
\noindent\hangindent=25bp\hangafter=1\hspace*{6bp}%
\makebox[19bp][r]{\hypertarget{bib-mumford-4}{}4.\enspace}Varieties defined by quadratic equations (with an Appendix by G. Kempf), in \emph{Questions on Algebraic Varieties}, Centro Internationale Matematica Estivo, Cremonese, Rome (1970), 29--100.\par
\noindent\hangindent=25bp\hangafter=1\hspace*{6bp}%
\makebox[19bp][r]{\hypertarget{bib-480-6}{}\hypertarget{bib-480-7}{}\hypertarget{bib-mumford-5}{}5.\enspace}\emph{Abelian Varieties}, Oxford Univ. Press, Oxford (1970), ix + 242 pp.\par
\noindent\hangindent=25bp\hangafter=1\hspace*{6bp}%
\makebox[19bp][r]{\hypertarget{bib-mumford-6}{}6.\enspace}The topology of normal singularities of an algebraic surface and a criterion for simplicity, \emph{Publ. Math. IHES 9} (1961) 5--22.\par
\par\vspace{2.15pt}\noindent \hypertarget{bib-480-8}{}Nagata, M.\par
\noindent\hangindent=25bp\hangafter=1\hspace*{6bp}%
\makebox[19bp][r]{\hypertarget{bib-480-9}{}1.\enspace}On the embedding problem of abstract varieties in projective varieties, \emph{Mem. Coll. Sci. Kyoto (A) 30} (1956), 71--82.\par
\noindent\hangindent=25bp\hangafter=1\hspace*{6bp}%
\makebox[19bp][r]{\hypertarget{bib-480-10}{}\hypertarget{bib-nagata-2}{}2.\enspace}A general theory of algebraic geometry over Dedekind domains, \emph{Amer. J. of Math. 78} (1956), 78--116.\par
\noindent\hangindent=25bp\hangafter=1\hspace*{6bp}%
\makebox[19bp][r]{\hypertarget{bib-480-11}{}3.\enspace}On the imbeddings of abstract surfaces in projective varieties, \emph{Mem. Coll. Sci. Kyoto (A) 30} (1957), 231--235.\par
\noindent\hangindent=25bp\hangafter=1\hspace*{6bp}%
\makebox[19bp][r]{\hypertarget{bib-480-12}{}4.\enspace}Existence theorems for non-projective complete algebraic varieties, \emph{Ill. J. Math. 2} (1958), 490--498.\par
\noindent\hangindent=25bp\hangafter=1\hspace*{6bp}%
\makebox[19bp][r]{\hypertarget{bib-nagata-5}{}5.\enspace}On rational surfaces I, II, \emph{Mem. Coll. Sci. Kyoto (A) 32} (1960), 351--370, and \emph{33} (1960), 271--293.\par
\noindent\hangindent=25bp\hangafter=1\hspace*{6bp}%
\makebox[19bp][r]{\hypertarget{bib-480-13}{}\hypertarget{bib-nagata-6}{}6.\enspace}Imbedding of an abstract variety in a complete variety, \emph{J. Math. Kyoto Univ. 2} (1962), 1--10.\par
\noindent\hangindent=25bp\hangafter=1\hspace*{6bp}%
\makebox[19bp][r]{\hypertarget{bib-480-14}{}\hypertarget{bib-480-15}{}7.\enspace}\emph{Local Rings}, Interscience Tracts in Pure \& Applied Math. \emph{13}, J. Wiley, New York (1962).\par
\noindent\hangindent=25bp\hangafter=1\hspace*{6bp}%
\makebox[19bp][r]{\hypertarget{bib-nagata-8}{}8.\enspace}On self-intersection number of a section on a ruled surface, \emph{Nagoya Math. J. 37} (1970), 191--196.\par
\par\vspace{2.15pt}\noindent \hypertarget{bib-480-16}{}Nakai, Y.\par
\noindent\hangindent=25bp\hangafter=1\hspace*{6bp}%
\makebox[19bp][r]{\hypertarget{bib-480-17}{}1.\enspace}A criterion of an ample sheaf on a projective scheme, \emph{Amer. J. Math. 85} (1963), 14--26.\par
\noindent\hangindent=25bp\hangafter=1\hspace*{6bp}%
\makebox[19bp][r]{\hypertarget{bib-nakai-2}{}2.\enspace}Some fundamental lemmas on projective schemes, \emph{Trans. Amer. Math. Soc. 109} (1963), 296--302.\par
\par\vspace{2.15pt}\noindent \hypertarget{bib-480-18}{}Nakano, S.\par
\noindent\hangindent=25bp\hangafter=1\hspace*{6bp}%
\makebox[19bp][r]{\hypertarget{bib-nakano-1}{}1.\enspace}On complex analytic vector bundles, \emph{J. Math. Soc. Japan, 7} (1955) 1--12.\par
\par\vspace{2.15pt}\noindent \hypertarget{bib-480-19}{}Narasimhan, M. S. and Seshadri, C. S.\par
\noindent\hangindent=25bp\hangafter=1\hspace*{6bp}%
\makebox[19bp][r]{\hypertarget{bib-narasimhan-seshadri-1}{}1.\enspace}Stable and unitary vector bundles on a compact Riemann surface, \emph{Annals of Math. 82} (1965), 540--567.\par
\par\vspace{2.15pt}\noindent \hypertarget{bib-480-20}{}Noether, M.\par
\noindent\hangindent=25bp\hangafter=1\hspace*{6bp}%
\makebox[19bp][r]{\hypertarget{bib-480-20-item-1}{}1.\enspace}\emph{Zur Grundlegung der Theorie der Algebraischen Raumcurven}, Verlag der K\"oniglichen Akademie der Wissenschaften, Berlin (1883).\par
\par\vspace{2.15pt}\noindent \hypertarget{bib-480-21}{}Olson, L.\par
\noindent\hangindent=25bp\hangafter=1\hspace*{6bp}%
\makebox[19bp][r]{\hypertarget{bib-olson-1}{}1.\enspace}An elementary proof that elliptic curves are abelian varieties, \emph{Ens. Math. 19} (1973), 173--181.\par
}
% Contents from body/page-481.tex.
% Original printed page 467.
\ChapterSheet[10.25]{481}{467}{\Goto{p-473}{Bibliography}}{%
\fontsize{9.25}{9.8}\selectfont%
\setlength{\parskip}{0pt}\setlength{\parindent}{0pt}%
\setlength{\emergencystretch}{1.7em}%
\par\vspace{2.15pt}\noindent Ramanujam, C. P.\par
\noindent\hangindent=25bp\hangafter=1\hspace*{6bp}%
\makebox[19bp][r]{\hypertarget{bib-ramanujam-1}{}1.\enspace}Remarks on the Kodaira vanishing theorem, \emph{J. Indian Math. Soc. (N.S.) 36} (1972), 41--51.\par
\par\vspace{2.15pt}\noindent Ramis, J. P. and Ruget, G.\par
\noindent\hangindent=25bp\hangafter=1\hspace*{6bp}%
\makebox[19bp][r]{\hypertarget{bib-ramis-ruget-1}{}1.\enspace}Complexes dualisants et th\'eor\`eme de dualit\'e en g\'eom\'etrie analytique complexe, \emph{Pub. Math. IHES 38} (1970), 77--91.\par
\par\vspace{2.15pt}\noindent Ramis, J. P., Ruget, G., and Verdier, J. L.\par
\noindent\hangindent=25bp\hangafter=1\hspace*{6bp}%
\makebox[19bp][r]{\hypertarget{bib-ramis-ruget-verdier-1}{}1.\enspace}Dualit\'e relative en g\'eom\'etrie analytique complexe, \emph{Invent. Math. 13} (1971), 261--283.\par
\par\vspace{2.15pt}\noindent Roberts, J.\par
\noindent\hangindent=25bp\hangafter=1\hspace*{6bp}%
\makebox[19bp][r]{\hypertarget{bib-roberts-1}{}1.\enspace}Chow's moving lemma, in \emph{Algebraic Geometry, Oslo 1970} (F. Oort, ed.), Wolters-Noordhoff (1972), 89--96.\par
\par\vspace{2.15pt}\noindent Roquette, P.\par
\noindent\hangindent=25bp\hangafter=1\hspace*{6bp}%
\makebox[19bp][r]{\hypertarget{bib-roquette-1}{}1.\enspace}Absch\"atzung der Automorphismenanzahl von Funktionenk\"orpern bei Primzahlcharakteristik, \emph{Math. Zeit. 117} (1970) 157--163.\par
\par\vspace{2.15pt}\noindent Rotman, J. J.\par
\noindent\hangindent=25bp\hangafter=1\hspace*{6bp}%
\makebox[19bp][r]{\hypertarget{bib-rotman-1}{}1.\enspace}\emph{Notes on Homological Algebra}, Van Nostrand Reinhold Math. Studies \emph{26}, New York (1970).\par
\par\vspace{2.15pt}\noindent Saint-Donat, B.\par
\noindent\hangindent=25bp\hangafter=1\hspace*{6bp}%
\makebox[19bp][r]{\hypertarget{bib-saint-donat-1}{}1.\enspace}On Petri's analysis of the linear system of quadrics through a canonical curve, \emph{Math. Ann. 206} (1973), 157--175.\par
\par\vspace{2.15pt}\noindent Sally, J.\par
\noindent\hangindent=25bp\hangafter=1\hspace*{6bp}%
\makebox[19bp][r]{\hypertarget{bib-sally-1}{}1.\enspace}Regular overrings of regular local rings, \emph{Trans. Amer. Math. Soc. 171} (1972) 291--300.\par
\par\vspace{2.15pt}\noindent Samuel, P.\par
\noindent\hangindent=25bp\hangafter=1\hspace*{6bp}%
\makebox[19bp][r]{\hypertarget{bib-samuel-1}{}1.\enspace}\emph{M\'ethodes d'Alg\`ebre Abstraite en G\'eom\'etrie Alg\'ebrique}, Ergebnisse \emph{4}, Springer-Verlag, Heidelberg (1955).\par
\noindent\hangindent=25bp\hangafter=1\hspace*{6bp}%
\makebox[19bp][r]{\hypertarget{bib-481-0}{}\hypertarget{bib-samuel-2}{}2.\enspace}\emph{Lectures on old and new results on algebraic curves} (notes by S. Anantharaman), Tata Inst. Fund. Res. (1966), 127 pp.\par
\noindent\hangindent=25bp\hangafter=1\hspace*{6bp}%
\makebox[19bp][r]{\hypertarget{bib-481-1}{}3.\enspace}\emph{Anneaux Factoriels} (r\'edaction de A. Micali), Soc. Mat. de S\~ao Paulo (1963) 97 pp.\par
\par\vspace{2.15pt}\noindent Schlessinger, M.\par
\noindent\hangindent=25bp\hangafter=1\hspace*{6bp}%
\makebox[19bp][r]{\hypertarget{bib-schlessinger-1}{}1.\enspace}Functors of Artin rings, \emph{Trans. Amer. Math. Soc. 130} (1968), 208--222.\par
\par\vspace{2.15pt}\noindent Serre, J.-P.\par
\noindent\hangindent=25bp\hangafter=1\hspace*{6bp}%
\makebox[19bp][r]{\hypertarget{bib-481-2}{}1.\enspace}Cohomologie et g\'eom\'etrie alg\'ebrique, \emph{Proc. ICM} (1954), vol.~III, 515--520.\par
\noindent\hangindent=25bp\hangafter=1\hspace*{6bp}%
\makebox[19bp][r]{\hypertarget{bib-481-3}{}2.\enspace}Un th\'eor\`eme de dualit\'e, \emph{Comm. Math. Helv. 29} (1955), 9--26.\par
\noindent\hangindent=25bp\hangafter=1\hspace*{6bp}%
\makebox[19bp][r]{\hypertarget{bib-481-4}{}\hypertarget{bib-serre-3}{}3.\enspace}Faisceaux alg\'ebriques coh\'erents, \emph{Ann. of Math. 61} (1955), 197--278.\par
\noindent\hangindent=25bp\hangafter=1\hspace*{6bp}%
\makebox[19bp][r]{\hypertarget{bib-481-5}{}\hypertarget{bib-serre-4}{}4.\enspace}G\'eom\'etrie alg\'ebrique et g\'eom\'etrie analytique, \emph{Ann. Inst. Fourier 6} (1956), 1--42.\par
\noindent\hangindent=25bp\hangafter=1\hspace*{6bp}%
\makebox[19bp][r]{\hypertarget{bib-serre-5}{}5.\enspace}Sur la cohomologie des vari\'et\'es alg\'ebriques, \emph{J. de Maths. Pures et Appl. 36} (1957), 1--16.\par
\noindent\hangindent=25bp\hangafter=1\hspace*{6bp}%
\makebox[19bp][r]{\hypertarget{bib-481-6}{}\hypertarget{bib-serre-6}{}6.\enspace}Sur la topologie des vari\'et\'es alg\'ebriques en caract\'eristique $p$, \emph{Symposium Int. de Topologia Algebraica, Mexico} (1958), 24--53.\par
\noindent\hangindent=25bp\hangafter=1\hspace*{6bp}%
\makebox[19bp][r]{\hypertarget{bib-481-7}{}\hypertarget{bib-481-8}{}\hypertarget{bib-serre-7}{}7.\enspace}\emph{Groupes Alg\'ebriques et Corps de Classes}, Hermann, Paris (1959).\par
\noindent\hangindent=25bp\hangafter=1\hspace*{6bp}%
\makebox[19bp][r]{\hypertarget{bib-serre-8}{}8.\enspace}Rationalit\'e des fonctions $\zeta$ des vari\'et\'es alg\'ebriques (d'apr\`es B. Dwork) \emph{S\'eminaire Bourbaki 198} (1960).\par
\noindent\hangindent=25bp\hangafter=1\hspace*{6bp}%
\makebox[19bp][r]{\hypertarget{bib-481-9}{}9.\enspace}Analogues k\"ahl\'eriens de certaines conjectures de Weil, \emph{Annals of Math. 71} (1960), 392--394.\par
\noindent\hangindent=25bp\hangafter=1\hspace*{6bp}%
\makebox[19bp][r]{\hypertarget{bib-481-10}{}10.\enspace}Zeta and $L$ functions, in \emph{Arithmetical Algebraic Geometry} (Schilling, ed.), Harper \& Row, New York (1965), 82--92.\par
\noindent\hangindent=25bp\hangafter=1\hspace*{6bp}%
\makebox[19bp][r]{\hypertarget{bib-481-11}{}11.\enspace}\emph{Alg\`ebre Locale---Multiplicit\'es} (r\'edig\'e par P. Gabriel), Lectures Notes in Math. \emph{11}, Springer-Verlag, Heidelberg (1965).\par
}
% Contents from body/page-482.tex.
% Original printed page 468.
\ChapterSheet[10.25]{482}{468}{\Goto{p-473}{Bibliography}}{%
\fontsize{9.25}{9.71}\selectfont%
\setlength{\parskip}{0pt}\setlength{\parindent}{0pt}%
\setlength{\emergencystretch}{1.7em}%
\noindent\hangindent=25bp\hangafter=1\hspace*{6bp}%
\makebox[19bp][r]{\hypertarget{bib-482-0}{}\hypertarget{bib-482-1}{}12.\enspace}Prolongement de faisceaux analytiques coh\'erents, \emph{Ann. Inst. Fourier 16} (1966), 363--374.\par
\noindent\hangindent=25bp\hangafter=1\hspace*{6bp}%
\makebox[19bp][r]{\hypertarget{bib-482-2}{}13.\enspace}Crit\`ere de rationalit\'e pour les surfaces alg\'ebriques (d'apr\`es K. Kodaira), \emph{S\'eminaire Bourbaki 146} (1957).\par
\noindent\hangindent=25bp\hangafter=1\hspace*{6bp}%
\makebox[19bp][r]{\hypertarget{bib-serre-14}{}14.\enspace}\emph{A Course in Arithmetic}, Graduate Texts in Math. \emph{7}, Springer-Verlag, Heidelberg (1973) 115 pp.\par
\par\vspace{2.15pt}\noindent \hypertarget{bib-482-3}{}Severi, F.\par
\noindent\hangindent=25bp\hangafter=1\hspace*{6bp}%
\makebox[19bp][r]{\hypertarget{bib-severi-1}{}1.\enspace}Intorno ai punti doppi impropri di una superficie generale dello spazio a quattro dimensioni, e a suoi punti tripli apparenti, \emph{Rend. Circ. Matem. Palermo 15} (1901), 33--51.\par
\noindent\hangindent=25bp\hangafter=1\hspace*{6bp}%
\makebox[19bp][r]{\hypertarget{bib-482-4}{}\hypertarget{bib-482-5}{}\hypertarget{bib-severi-2}{}2.\enspace}\emph{Vorlesungen \"uber Algebraische Geometrie} (transl. by E. L\"offler), Johnson Pub. (rpt., 1968; 1st. ed., Leipzig 1921).\par
\noindent\hangindent=25bp\hangafter=1\hspace*{6bp}%
\makebox[19bp][r]{\hypertarget{bib-severi-3}{}3.\enspace}\"Uber die Grundlagen der algebraischen Geometrie, \emph{Hamb. Abh. 9} (1933) 335--364.\par
\par\vspace{2.15pt}\noindent \hypertarget{bib-482-6}{}Shafarevich, I. R.\par
\noindent\hangindent=25bp\hangafter=1\hspace*{6bp}%
\makebox[19bp][r]{\hypertarget{bib-482-7}{}1.\enspace}\emph{Algebraic surfaces}, Proc. Steklov Inst. Math. \emph{75} (1965) (trans. by A.M.S. 1967).\par
\noindent\hangindent=25bp\hangafter=1\hspace*{6bp}%
\makebox[19bp][r]{2.\enspace}\emph{Basic Algebraic Geometry}, Grundlehren \emph{213}, Springer-Verlag, Heidelberg (1974), xv + 439 pp.\par
\par\vspace{2.15pt}\noindent \hypertarget{bib-482-8}{}Shannon, D. L.\par
\noindent\hangindent=25bp\hangafter=1\hspace*{6bp}%
\makebox[19bp][r]{\hypertarget{bib-482-8-item-1}{}1.\enspace}Monoidal transforms of regular local rings, \emph{Amer. J. Math. 95} (1973) 294--320.\par
\par\vspace{2.15pt}\noindent \hypertarget{bib-482-9}{}Shioda, T.\par
\noindent\hangindent=25bp\hangafter=1\hspace*{6bp}%
\makebox[19bp][r]{\hypertarget{bib-482-9-item-1}{}1.\enspace}An example of unirational surfaces in characteristic $p$, \emph{Math. Ann. 211} (1974) 233--236.\par
\par\vspace{2.15pt}\noindent Siegel, C. L.\par
\noindent\hangindent=25bp\hangafter=1\hspace*{6bp}%
\makebox[19bp][r]{\hypertarget{bib-siegel-1}{}1.\enspace}Meromorphe Funktionen auf kompakten analytischen Mannigfaltigkeiten, \emph{Nach. Akad. Wiss. G\"ottingen} (1955), 71--77.\par
\par\vspace{2.15pt}\noindent \hypertarget{bib-482-10}{}Singh, B.\par
\noindent\hangindent=25bp\hangafter=1\hspace*{6bp}%
\makebox[19bp][r]{\hypertarget{bib-singh-1}{}1.\enspace}On the group of automorphisms of a function field of genus at least two, \emph{J. Pure Appl. Math. 4} (1974) 205--229.\par
\par\vspace{2.15pt}\noindent \hypertarget{bib-482-11}{}Spanier, E. H.\par
\noindent\hangindent=25bp\hangafter=1\hspace*{6bp}%
\makebox[19bp][r]{\hypertarget{bib-spanier-1}{}1.\enspace}\emph{Algebraic Topology}, McGraw-Hill, New York (1966).\par
\par\vspace{2.15pt}\noindent \hypertarget{bib-482-12}{}Stichtenoth, H.\par
\noindent\hangindent=25bp\hangafter=1\hspace*{6bp}%
\makebox[19bp][r]{\hypertarget{bib-stichtenoth-1}{}1.\enspace}\"Uber die Automorphismengruppe eines algebraischen Funktionenk\"orpers von Primzahlcharakteristik, \emph{Archiv der Math. 24} (1973) 527--544.\par
\par\vspace{2.15pt}\noindent \hypertarget{bib-482-13}{}Suominen, K.\par
\noindent\hangindent=25bp\hangafter=1\hspace*{6bp}%
\makebox[19bp][r]{\hypertarget{bib-suominen-1}{}1.\enspace}Duality for coherent sheaves on analytic manifolds, \emph{Ann. Acad. Sci. Fenn (A) 424} (1968), 1--19.\par
\par\vspace{2.15pt}\noindent \hypertarget{bib-482-14}{}Tate, J. T.\par
\noindent\hangindent=25bp\hangafter=1\hspace*{6bp}%
\makebox[19bp][r]{\hypertarget{bib-482-15}{}1.\enspace}Algebraic cycles and poles of the zeta function, in \emph{Arithmetical Algebraic Geometry} (Schilling, ed.), Harper \& Row, New York (1965), 93--110.\par
\noindent\hangindent=25bp\hangafter=1\hspace*{6bp}%
\makebox[19bp][r]{\hypertarget{bib-tate-2}{}\hypertarget{bib-482-16}{}2.\enspace}Residues of differentials on curves, \emph{Ann. Sci. de l'E.N.S.} (4) \emph{1} (1968), 149--159.\par
\noindent\hangindent=25bp\hangafter=1\hspace*{6bp}%
\makebox[19bp][r]{\hypertarget{bib-tate-3}{}3.\enspace}The arithmetic of elliptic curves, \emph{Inv. Math. 23} (1974), 179--206.\par
\par\vspace{2.15pt}\noindent \hypertarget{bib-482-17}{}Tjurin, A. N.\par
\noindent\hangindent=25bp\hangafter=1\hspace*{6bp}%
\makebox[19bp][r]{\hypertarget{bib-tjurin-1}{}1.\enspace}On the classification of two-dimensional fibre bundles over an algebraic curve of arbitrary genus (in Russian) \emph{Izv. Akad. Nauk SSSR Ser. Mat. 28} (1964) 21--52; \textbf{MR} \emph{29} (1965) \# 4762.\par
\noindent\hangindent=25bp\hangafter=1\hspace*{6bp}%
\makebox[19bp][r]{\hypertarget{bib-482-18}{}2.\enspace}Classification of vector bundles over an algebraic curve of arbitrary genus, \emph{Amer. Math. Soc. Translations 63} (1967) 245--279.\par
\par\vspace{2.15pt}\noindent \hypertarget{bib-482-19}{}Verdier, J.-L.\par
\noindent\hangindent=25bp\hangafter=1\hspace*{6bp}%
\makebox[19bp][r]{\hypertarget{bib-482-19-item-1}{}1.\enspace}Base change for twisted inverse image of coherent sheaves, in \emph{Algebraic Geometry, Bombay 1968}, Oxford Univ. Press, Oxford (1969), 393--408.\par
}
% Contents from body/page-483.tex.
% Original printed page 469.
\ChapterSheet[10.25]{483}{469}{\Goto{p-473}{Bibliography}}{%
\fontsize{9.25}{9.8}\selectfont%
\setlength{\parskip}{0pt}\setlength{\parindent}{0pt}%
\setlength{\emergencystretch}{1.7em}%
\par\vspace{2.15pt}\noindent \hypertarget{bib-483-0}{}\hypertarget{bib-vitushkin-1}{}Vitushkin, A. G.\par
\noindent\hangindent=25bp\hangafter=1\hspace*{6bp}%
\makebox[19bp][r]{\hypertarget{bib-vitushkin-1-item-1}{}1.\enspace}On polynomial transformation of $\mathbf C^n$, in \emph{Manifolds, Tokyo 1973}, Tokyo Univ. Press, Tokyo (1975), 415--417.\par
\par\vspace{2.15pt}\noindent \hypertarget{bib-483-1}{}Walker, R. J.\par
\noindent\hangindent=25bp\hangafter=1\hspace*{6bp}%
\makebox[19bp][r]{1.\enspace}\emph{Algebraic Curves}, Princeton Univ., Princeton (1950), Dover reprint (1962).\par
\par\vspace{2.15pt}\noindent \hypertarget{bib-483-2}{}van der Waerden, B. L.\par
\noindent\hangindent=25bp\hangafter=1\hspace*{6bp}%
\makebox[19bp][r]{\hypertarget{bib-483-2-item-1}{}1.\enspace}\emph{Modern Algebra}, Frederick Ungar Pub. Co, New York: I (1953), xii + 264 pp.; II (1950), ix + 222 pp.\par
\par\vspace{2.15pt}\noindent \hypertarget{bib-483-3}{}\hypertarget{bib-weil-1}{}Weil, A.\par
\noindent\hangindent=25bp\hangafter=1\hspace*{6bp}%
\makebox[19bp][r]{\hypertarget{bib-483-4}{}1.\enspace}\emph{Foundations of Algebraic Geometry}, Amer. Math. Soc., Colloquium Publ. \emph{29} (1946) (revised and enlarged edition 1962), xx + 363 pp.\par
\noindent\hangindent=25bp\hangafter=1\hspace*{6bp}%
\makebox[19bp][r]{\hypertarget{bib-weil-2}{}2.\enspace}\emph{Sur les Courbes Alg\'ebriques et les Vari\'et\'es qui s'en D\'eduisent}, Hermann, Paris (1948). This volume and the next have been republished in one volume, \emph{Courbes Alg\'ebriques et Vari\'et\'es Ab\'eliennes}, Hermann, Paris (1971), 249 pp.\par
\noindent\hangindent=25bp\hangafter=1\hspace*{6bp}%
\makebox[19bp][r]{\hypertarget{bib-483-5}{}\hypertarget{bib-483-6}{}3.\enspace}\emph{Vari\'et\'es Ab\'eliennes et Courbes Alg\'ebriques}, Hermann, Paris (1948).\par
\noindent\hangindent=25bp\hangafter=1\hspace*{6bp}%
\makebox[19bp][r]{\hypertarget{bib-483-7}{}4.\enspace}Number of solutions of equations over finite fields, \emph{Bull. Amer. Math. Soc. 55} (1949), 497--508.\par
\noindent\hangindent=25bp\hangafter=1\hspace*{6bp}%
\makebox[19bp][r]{\hypertarget{bib-483-8}{}5.\enspace}\emph{Vari\'et\'es K\"ahl\'eriennes}, Hermann, Paris (1958), 175 pp.\par
\noindent\hangindent=25bp\hangafter=1\hspace*{6bp}%
\makebox[19bp][r]{\hypertarget{bib-weil-6}{}6.\enspace}On the projective embedding of abelian varieties, in \emph{Algebraic Geometry and Topology (in honor of S. Lefschetz)}, Princeton Univ., Princeton (1957) 177--181.\par
\par\vspace{2.15pt}\noindent \hypertarget{bib-483-9}{}Wells, R. O., Jr.\par
\noindent\hangindent=25bp\hangafter=1\hspace*{6bp}%
\makebox[19bp][r]{1.\enspace}\emph{Differential Analysis on Complex Manifolds}, Prentice-Hall (1973), x + 252 pp.\par
\par\vspace{2.15pt}\noindent \hypertarget{bib-483-10}{}Weyl, H.\par
\noindent\hangindent=25bp\hangafter=1\hspace*{6bp}%
\makebox[19bp][r]{\hypertarget{bib-weyl-1}{}1.\enspace}\emph{Die Idee der Riemannschen Fl\"ache}, Teubner (3rd ed., 1955), vii + 162 pp. (1st ed., 1913).\par
\par\vspace{2.15pt}\noindent \hypertarget{bib-483-11}{}Zariski, O.\par
\noindent\hangindent=25bp\hangafter=1\hspace*{6bp}%
\makebox[19bp][r]{\hypertarget{bib-483-12}{}1.\enspace}The concept of a simple point on an abstract algebraic variety, \emph{Trans. Amer. Math. Soc. 62} (1947), 1--52.\par
\noindent\hangindent=25bp\hangafter=1\hspace*{6bp}%
\makebox[19bp][r]{\hypertarget{bib-483-13}{}\hypertarget{bib-zariski-2}{}2.\enspace}A simple analytical proof of a fundamental property of birational transformations, \emph{Proc. Nat. Acad. Sci. USA 35} (1949), 62--66.\par
\noindent\hangindent=25bp\hangafter=1\hspace*{6bp}%
\makebox[19bp][r]{\hypertarget{bib-483-14}{}3.\enspace}\emph{Theory and Applications of Holomorphic Functions on Algebraic Varieties over Arbitrary Ground Fields}, Memoirs of Amer. Math. Soc. New York (1951).\par
\noindent\hangindent=25bp\hangafter=1\hspace*{6bp}%
\makebox[19bp][r]{\hypertarget{bib-zariski-4}{}4.\enspace}Complete linear systems on normal varieties and a generalization of a lemma of Enriques-Severi, \emph{Ann. of Math. 55} (1952), 552--592.\par
\noindent\hangindent=25bp\hangafter=1\hspace*{6bp}%
\makebox[19bp][r]{\hypertarget{bib-483-15}{}\hypertarget{bib-483-16}{}\hypertarget{bib-zariski-5}{}5.\enspace}\emph{Introduction to the Problem of Minimal Models in the Theory of Algebraic Surfaces}, Pub. Math. Soc. of Japan \emph{4} (1958), vii + 89 pp.\par
\noindent\hangindent=25bp\hangafter=1\hspace*{6bp}%
\makebox[19bp][r]{\hypertarget{bib-483-17}{}\hypertarget{bib-zariski-6}{}6.\enspace}The problem of minimal models in the theory of algebraic surfaces, \emph{Amer. J. Math. 80} (1958), 146--184.\par
\noindent\hangindent=25bp\hangafter=1\hspace*{6bp}%
\makebox[19bp][r]{\hypertarget{bib-483-18}{}7.\enspace}The theorem of Riemann-Roch for high multiples of an effective divisor on an algebraic surface, \emph{Ann. Math. 76} (1962), 560--615.\par
\noindent\hangindent=25bp\hangafter=1\hspace*{6bp}%
\makebox[19bp][r]{\hypertarget{bib-zariski-8}{}8.\enspace}\emph{Collected papers. Vol. I. Foundations of Algebraic Geometry and Resolution of Singularities}, ed. H. Hironaka and D. Mumford, M.I.T. Press, Cambridge (1972), xxi + 543 pp. Vol.~II, \emph{Holomorphic Functions and Linear Systems}, ed. M. Artin and D. Mumford, M.I.T. Press (1973), xxiii + 615 pp.\par
\noindent\hangindent=25bp\hangafter=1\hspace*{6bp}%
\makebox[19bp][r]{\hypertarget{bib-483-19}{}\hypertarget{bib-483-20}{}9.\enspace}On Castelnuovo's criterion of rationality $p_a=P_2=0$ of an algebraic surface, \emph{Ill. J. Math. 2} (1958) 303--315.\par
\noindent\hangindent=25bp\hangafter=1\hspace*{6bp}%
\makebox[19bp][r]{\hypertarget{bib-zariski-10}{}10.\enspace}\emph{Algebraic Surfaces}, 2nd suppl. ed., Ergebnisse \emph{61}, Springer-Verlag, Heidelberg (1971).\par
\par\vspace{2.15pt}\noindent \hypertarget{bib-483-21}{}Zariski, O. and Samuel, P.\par
\noindent\hangindent=25bp\hangafter=1\hspace*{6bp}%
\makebox[19bp][r]{\hypertarget{bib-zariski-samuel-1}{}\hypertarget{legacy-b-zs}{}1.\enspace}\emph{Commutative Algebra (Vol. I, II)}, Van Nostrand, Princeton (1958, 1960).\par
}
